\section{Metodologia}

\subsection{Delineamento, catálogo e relevância}

Realizamos um estudo experimental de recuperação textual, seguido de avaliações complementares e validação funcional do protótipo. Utilizamos uma cópia local do banco disponibilizado para o projeto. Dos 873 registros de logradouros, excluímos cinco sem nome e 15 com marcadores genéricos, como ``nao oficial'' e ``n.o''. O índice resultante contém 853 registros e 847 nomes normalizados. Trata-se do recorte disponível nesse backup, sem pressupor cobertura integral do acervo Pauliceia, dos séculos XIX e XX ou dos livros de emplacamento mencionados no planejamento.

A normalização para agrupamento e cálculo das similaridades aplicou decomposição Unicode NFD, remoção de marcas de acentuação, conversão por \texttt{casefold} e regularização dos espaços. Mantivemos os nomes originais para exibição e não modernizamos suas grafias no cadastro. Cada consulta recebeu como relevantes todos os IDs com o mesmo nome normalizado de sua referência. Essa decisão avalia a correspondência textual: registros homônimos podem representar situações espaciais diferentes, que o nome isolado não resolve.

\subsection{Consultas e separação dos dados}

O conjunto principal reuniu 6.136 consultas: 847 nomes do cadastro e 5.289 variantes sintéticas. Geramos variantes de normalização, abreviação do tipo de logradouro, omissão, substituição, inserção e transposição de caracteres, duas omissões em palavras distintas e inversão da ordem das palavras após o primeiro termo. As posições de edição foram determinadas por hash e semente fixa. A substituição utilizou a letra seguinte no alfabeto e a inserção repetiu um caractere. Portanto, esses exemplos representam perturbações controladas; não reproduzem frequências observadas de erros de teclado, OCR ou uso real.

As regras tiveram condições de aplicação, de modo que nem todo nome recebeu todas as variantes. O gerador removeu uma repetição e verificou colisões entre consultas e referências distintas; não houve exclusões por ambiguidade nessa execução. As nove categorias totalizaram 847 originais, 22 normalizações, 834 abreviações, 834 omissões, 834 substituições, 834 inserções, 834 transposições, 491 casos de duas omissões e 606 inversões de ordem.

Para reduzir a transferência de variantes relacionadas entre partições, agrupamos nomes após normalização, desconsiderando a ordem das palavras e os sufixos de lado par/ímpar. Reservamos os seis grupos utilizados no piloto para desenvolvimento. Os demais foram ordenados por hash com a semente \texttt{pauliceia-tcc-v1} e distribuídos aproximadamente em 60\% para treino, 20\% para validação e 20\% para teste. A Tabela~\ref{tab:particoes} apresenta as quantidades efetivas. Consultas do conjunto de desenvolvimento não participaram do ajuste nem da seleção do modelo.

% Gerado por gerar_tabelas.py a partir dos resultados congelados.
\begin{table}[tb]
\caption{Separação do conjunto principal por grupos de nomes.}
\label{tab:particoes}
\centering
\small
\setlength{\tabcolsep}{3pt}
\begin{tabular}{lrrr}
\toprule
Partição & Grupos & Nomes & Consultas \\
\midrule
Desenvolvimento & 6 & 6 & 43 \\
Treino & 496 & 505 & 3681 \\
Validacao & 166 & 169 & 1205 \\
Teste & 166 & 167 & 1207 \\
\bottomrule
\end{tabular}

\end{table}


O índice manteve o catálogo completo em todas as etapas, porque recuperar um nome exige que ele esteja indexado. A separação incidiu sobre os grupos de consultas e seus rótulos para ajuste e seleção, não sobre a disponibilidade dos documentos candidatos. Consequentemente, nomes associados às consultas de teste podem aparecer como candidatos negativos em consultas de treino. O classificador recebeu medidas numéricas de similaridade, sem IDs ou nomes como atributos categóricos. O protocolo avalia generalização para novos grupos de consultas dentro de um catálogo fixo, não para outro catálogo.

\subsection{Recuperação de candidatos}

Configuramos o Elasticsearch 8.19.4 com um fragmento primário e sem réplicas. O campo de palavras utilizou tokenização \texttt{standard}, conversão para minúsculas e \texttt{asciifolding}. Um segundo campo acrescentou um filtro de n-gramas com tamanho mínimo e máximo iguais a três, aplicado aos tokens. Ambos usaram BM25 com os parâmetros padrão, $k_1=1{,}2$ e $b=0{,}75$~\cite{elasticSimilarity}.

A referência BM25 consultou apenas o campo de palavras. A alternativa BM25 com trigramas utilizou \texttt{multi\_match} do tipo \texttt{most\_fields}, combinando os escores dos dois campos sem pesos adicionais~\cite{elasticMulti}. Cada método recuperou até 100 candidatos, ordenados por escore decrescente e ID crescente no desempate. O limite foi mantido nos experimentos e no serviço, sem busca posterior por valores alternativos. BM25 e trigramas podem produzir conjuntos distintos; todos os reordenadores operaram sobre a mesma lista obtida pelos trigramas.

\subsection{Similaridades e reordenação supervisionada}

Representamos cada par consulta--candidato por sete atributos: similaridades normalizadas de Damerau--Levenshtein, Jaro--Winkler e Levenshtein; Token Sort e Token Set divididos por 100; escore do Elasticsearch dividido pelo maior escore da consulta; e razão entre o menor e o maior comprimento das strings. Usamos RapidFuzz 3.14.3 e matrizes \texttt{float32}, com as mesmas funções no experimento e no serviço online.

Damerau--Levenshtein considera operações de edição, incluindo transposições adjacentes; sua similaridade normalizada utiliza o maior comprimento como referência~\cite{rapidDL}. Token Sort ordena as palavras antes da comparação, enquanto Token Set considera palavras únicas e comuns. Essas funções empregam a similaridade Indel normalizada do RapidFuzz, e não devem ser descritas simplesmente como distância de Levenshtein sobre tokens~\cite{rapidFuzz}. As diferentes representações fornecem sinais que podem responder de maneira distinta a alterações locais e à ordem das palavras.

Treinamos um Random Forest, família de modelos baseada na combinação de árvores~\cite{breiman2001}, para classificar pares como relevantes ou não relevantes. A implementação \texttt{RandomForestClassifier}, do scikit-learn 1.7.2~\cite{sklearnRF}, recebeu 347.363 pares, dos quais 3.678 relevantes. Cada par teve peso inversamente proporcional ao número de candidatos de sua consulta, e as classes foram balanceadas em cada amostra bootstrap com \texttt{balanced\_subsample}. A diferença entre consultas e pares positivos decorre da quantidade de relevantes recuperados, que pode ser nula ou maior que um.

Fixamos 200 árvores e semente 42. A grade combinou profundidade máxima 12 ou ilimitada com mínimo de dois ou cinco exemplos por folha; os demais parâmetros permaneceram nos padrões da versão registrada. Selecionamos a configuração pelo maior MRR médio por nome na validação, com desempate pelo acerto na primeira posição por nome e, depois, pela ordem da grade. O modelo selecionado teve profundidade 12 e mínimo de dois exemplos por folha. O treinamento usou até quatro tarefas paralelas e a previsão uma tarefa, preservando a ordem de soma das árvores.

A saída da classe relevante em \texttt{predict\_proba} foi utilizada como escore de ordenação, sem assumir calibração probabilística. Empates conservaram a ordem dos trigramas. O modelo foi serializado, e as previsões antes e depois da leitura do arquivo foram comparadas. Após a seleção pela validação, o modelo permaneceu congelado para o teste e as avaliações adicionais.

Como comparadores adicionais, reordenamos os candidatos por Damerau--Levenshtein, Jaro--Winkler, Token Sort e Token Set individualmente, com o mesmo desempate. Essa comparação foi acrescentada após a observação do teste original, sem retreinamento, novos pesos ou escolha de hiperparâmetros. Ela compara métodos finais; não constitui remoção controlada de atributos do Random Forest, que recebe sete medidas.

\subsection{Métricas de recuperação}

Para uma consulta $q$, sejam $G_q$ o conjunto de IDs relevantes e $L_q^k$ os IDs nas primeiras $k$ posições. Calculamos o acerto em $k$ posições, a precisão e a cobertura:
\begin{align}
 H@k(q) &= \mathbf{1}[G_q\cap L_q^k\ne\varnothing],\\
 P@k(q) &= \frac{|G_q\cap L_q^k|}{k},\\
 R@100(q) &= \frac{|G_q\cap L_q^{100}|}{|G_q|}.
\end{align}
O recíproco da posição do primeiro relevante, limitado aos 100 candidatos, foi definido por
\begin{equation}
 RR@100(q)=\begin{cases}1/r_q,&r_q\leq100,\\0,&\text{sem relevante recuperado},\end{cases}
\end{equation}
em que $r_q$ é a primeira posição relevante. MRR@100 corresponde à média de $RR@100$. Reportamos Top-1 como a média de $H@1$, além de $H@3$, $H@5$, $P@1$, $P@3$, $P@5$ e Recall@100. Uma consulta sem alvo entre os candidatos continuou contribuindo com zero para acerto e MRR. A reordenação não introduziu IDs ausentes.

As tabelas principais usam médias por consulta. Também preservamos médias por nome e resultados por transformação. Acerto@K e Precision@K têm interpretações diferentes: encontrar uma referência entre cinco sugestões não torna todas as cinco relevantes. As comparações de ganhos e perdas são descritivas, sem testes de significância ou intervalos de confiança. As variantes de uma mesma referência não são observações independentes para inferência estatística.

\subsection{Avaliação complementar e fontes documentais}

Após o teste original, construímos um complemento exploratório de 179 consultas referentes a 168 nomes: 90 variantes históricas simuladas, 79 fonéticas simuladas e dez grafias documentadas. As simulações aplicaram substituições textuais controladas. O grupo chamado fonético não resulta de gravações de fala ou validação fonológica, e o grupo histórico simulado não corresponde a ocorrências documentais verificadas para cada transformação.

As dez consultas documentadas foram associadas a sete fontes municipais digitais. As ementas da Lei 1.734/1913 e do Ato 3.330/1930 registram, respectivamente, Higienópolis e Teodoro Sampaio~\cite{lei1734,ato3330}. Outras fontes registram Santa Ifigênia~\cite{santaIfigenia}, Artur Azevedo~\cite{conpresp7}, Bela Cintra, Barão de Tatuí, Humaitá e Triunfo~\cite{smsub40}, Travessa do Comércio~\cite{lei1552} e Beco dos Aflitos~\cite{aflitos}. O manifesto preservou URL, trecho localizador, consulta e ID associado.

A evidência documental atesta a ocorrência da grafia consultada. O vínculo dessa grafia ao ID do backup foi uma inferência textual sem revisão histórica independente. Ementas digitais de normas antigas podem refletir modernizações editoriais; a data da norma não foi tratada como comprovação da grafia no documento original. Não inferimos continuidade da extensão física da rua a partir dessas ocorrências.

O complemento incluiu nomes de diferentes partições e sete consultas já presentes nos conjuntos anteriores, incluindo o piloto. Identificamos separadamente 32 consultas novas ligadas a nomes do teste original. Nenhuma dessas avaliações foi usada para modificar o modelo. A construção posterior e a exposição anterior de parte dos nomes impedem tratar as 179 consultas como um novo teste independente.

\subsection{Integração e validação do portal}

O portal local passou a apresentar cinco sugestões com nome e ID após a recuperação de 100 candidatos e a reordenação. O usuário seleciona uma rua e informa número e ano; a API consulta os pontos cadastrados no PostgreSQL/PostGIS. O novo fluxo não interpola coordenadas. Os intervalos temporais são inclusivos; ano final nulo ou zero é interpretado como 2000, convenção herdada da API e não evidência histórica adicional.

A arquitetura usou Elasticsearch para recuperação e um serviço Python para similaridades e previsão. O planejamento inicial previa Painless no indexador; essa implementação não foi realizada. O serviço Python permitiu compartilhar funções, versões e representação numérica com o experimento, evitando uma segunda implementação cuja equivalência precisaria ser verificada. Não comparamos desempenho de Python com Painless. O serviço carrega o modelo uma vez, verifica sua assinatura e fornece identificação do modelo na resposta; indisponibilidade resulta em erro explícito.

A amostra de latência foi fixada antes da medição: 112 consultas, sendo 82 do teste e 30 do complemento. Após uma chamada de aquecimento por consulta, executamos três rodadas sequenciais, totalizando 336 buscas medidas. O cronômetro monotônico abrangeu a requisição HTTP pelo proxy local, o processamento e a leitura do JSON. A busca incluiu API, Elasticsearch, similaridades e previsão. Foram usadas novas conexões HTTP, sem concorrência. O P95 foi calculado por interpolação linear na posição $(n-1)\times0{,}95$ dos tempos ordenados.

Para a verificação espacial, um leitor independente da rota extraiu coordenadas e intervalos do backup. Dos 5.872 registros de pontos, 20 não tinham geometria ou ano inicial necessários, restando 5.852 no oráculo. Um registro com ano inicial 199 foi sinalizado e preservado, sem correção presumida; endereços associados a essa anomalia não serviram como âncoras da seleção temporal. Os 219 controles incluíram 103 respostas com coordenadas e 116 vazias, abrangendo 37 ruas e um ID inexistente. Houve nove casos com múltiplas coordenadas e 21 com numeração decimal. A comparação usou conjuntos de coordenadas, sem dependência da ordem, e tolerância de $10^{-9}$ grau por componente.

Também executamos três rodadas de 20 fluxos compostos por busca, seleção explícita do ID esperado entre as sugestões e consulta de pontos. Os 60 fluxos não representam 60 endereços independentes. A medição excluiu renderização do navegador, digitação e decisão humana. A verificação espacial avaliou fidelidade ao backup, sem levantamento externo para medir erro geográfico em metros.

\subsection{Ambiente e rastreabilidade}

O ambiente local utilizou contêineres Linux, PostgreSQL 16/PostGIS 3.5, Elasticsearch 8.19.4 e API Node.js 22, integrada à interface Vue/OpenLayers. O treinamento registrou Python 3.13.3, NumPy 2.2.5, scikit-learn 1.7.2 e RapidFuzz 3.14.3. O cliente da medição HTTP registrou Windows 11 e 16 processadores lógicos; modelo de CPU, memória e recursos atribuídos aos contêineres não foram registrados no protocolo. Os tempos devem ser interpretados nesse limite de caracterização do ambiente.

Os artefatos locais incluem consultas, partições, candidatos, escores, métricas individuais, parâmetros, versões e hashes das entradas e do modelo. As avaliações adicionais reutilizaram os candidatos salvos e verificaram a reprodução dos resultados anteriores. Essa rastreabilidade permite conferir as tabelas e repetir os procedimentos com os mesmos arquivos, sem implicar publicação de todo o banco de origem.
